% ====================================================================
% design.tex (design note example)
% ====================================================================

\section{Design}
\label{sec:design}

\subsection{Choice of circuit}

A single resistor and capacitor would be the simplest filter, but its response
falls at only \qty{20}{\decibel} per decade, so a corner frequency low enough
to give \qty{35}{\decibel} at \qty{9.5}{\kilo\hertz} would sit at around
\qty{170}{\hertz} and would cut into the signal band. A second-order filter
falls at \qty{40}{\decibel} per decade, which allows the corner to be placed at
\qty{1}{\kilo\hertz}, comfortably above the signal band, whilst still meeting
the stopband requirement.

The unity-gain Sallen-Key circuit~\cite{sallen_key_1955} in
\cref{fig:filterCircuit} was chosen because it needs only one op-amp, two
resistors and two capacitors, and because with the op-amp connected as a
follower the passband gain is exactly one and does not depend on any resistor
ratio. A Butterworth response is used, as it is the flattest in the passband
and requirement R4 is about flatness.

% ====================================================================
% filterCircuit.tex
% --------------------------------------------------------------------
% Unity-gain Sallen-Key low-pass filter, drawn with CircuiTikZ.
% ====================================================================

\begin{figure}[htbp]
\centering
\begin{circuitikz}
  % Op-amp as a follower, non-inverting input at the top
  \node[op amp, noinv input up] (U1) at (8,0) {U1};

  % Input, R1, junction A, R2, junction B, into the non-inverting input
  \draw (0,0 |- U1.+) node[left] {$\Vin$}
        to[short, o-] ++(0.4,0)
        to[R, l=$R_1$] ++(2.4,0) coordinate (A)
        to[R, l=$R_2$] ++(2.6,0) coordinate (B)
        -- (U1.+);

  % C2 from junction B to ground
  \draw (B) to[C, l=$C_2$, *-] ++(0,-2.2) node[ground] {};

  % C1 from junction A up and across to the output
  \draw (A) to[short, *-] ++(0,2) coordinate (Atop)
        to[C, l=$C_1$] (Atop -| U1.out)
        to[short, -*] (U1.out);

  % Follower feedback from the output to the inverting input
  \draw (U1.-) -- ++(-0.5,0) -- ++(0,-1.3) coordinate (fb)
        -- (fb -| U1.out) to[short, -*] ([yshift=-0.01cm]U1.out);

  % Output
  \draw (U1.out) to[short, -o] ++(1.2,0) node[right] {$\Vout$};
\end{circuitikz}
\caption{The unity-gain Sallen-Key low-pass filter. $C_1$ is the feedback
capacitor and $C_2$ is the capacitor to ground.}
\label{fig:filterCircuit}
\end{figure}


\subsection{Design equations}

The transfer function of the circuit, which is derived in
\cref{app:derivation}, is

\begin{equation}
  \label{eq:transferFunction}
  H(s) = \frac{\Vout}{\Vin}
       = \frac{1}{1 + s\,C_2(R_1 + R_2) + s^2 R_1 R_2 C_1 C_2}.
\end{equation}

Comparing this with the standard second-order form gives the natural frequency
and the quality factor,

\begin{align}
  \label{eq:naturalFrequency}
  f_0 &= \frac{1}{2\pi\sqrt{R_1 R_2 C_1 C_2}}, \\[4pt]
  \label{eq:qualityFactor}
  Q   &= \frac{\sqrt{R_1 R_2 C_1 C_2}}{C_2(R_1 + R_2)}.
\end{align}

If the two resistors are made equal, so that $R_1 = R_2 = R$,
\cref{eq:qualityFactor} reduces to $Q = \tfrac{1}{2}\sqrt{C_1/C_2}$, and the
quality factor then depends only on the ratio of the capacitors. A Butterworth
response needs $Q = 1/\sqrt{2}$, which means $C_1 = 2C_2$.

\subsection{Component values}

Capacitors come in far fewer values than resistors, so the capacitors are
chosen first. Taking $C_2$ as \qty{10}{\nano\farad}, $C_1$ has to be
\qty{20}{\nano\farad}, which is made from two \qty{10}{\nano\farad} parts in
parallel. This gives the exact two-to-one ratio, and it means all three
capacitors are the same component, so they will tend to drift together with
temperature and the ratio that sets $Q$ is preserved.

Rearranging \cref{eq:naturalFrequency} for a corner at \qty{1}{\kilo\hertz}
gives the resistor value,

\begin{equation}
  R = \frac{1}{2\pi f_0\sqrt{C_1 C_2}}
    = \frac{1}{2\pi \times \qty{1}{\kilo\hertz} \times \qty{14.14}{\nano\farad}}
    = \qty{11.25}{\kilo\ohm},
\end{equation}
and the nearest standard value is \qty{11.3}{\kilo\ohm}, which moves the
corner to \qty{996}{\hertz}. The full list of parts is in
\cref{tab:components}.

\begin{table}[htbp]
\centering
\begin{tabularx}{\textwidth}{@{} l l l X @{}}
\toprule
\textbf{Part} & \textbf{Value} & \textbf{Tolerance} & \textbf{Notes} \\
\midrule
$R_1$, $R_2$ & \qty{11.3}{\kilo\ohm} & \qty{1}{\percent} & Thin film, E96 series \\
$C_1$        & \qty{20}{\nano\farad} & \qty{5}{\percent} & Two \qty{10}{\nano\farad} C0G parts in parallel \\
$C_2$        & \qty{10}{\nano\farad} & \qty{5}{\percent} & C0G, the same part as $C_1$ \\
U1           & Op-amp                & n/a               & Rail-to-rail input and output, GBW of \qty{1}{\mega\hertz} or more \\
\bottomrule
\end{tabularx}
\caption{Component values for the filter.}
\label{tab:components}
\end{table}

C0G capacitors are specified because their value does not change with the
voltage across them, whereas the capacitance of the cheaper X7R type falls as
the voltage rises, which would distort the signal as well as moving the
corner frequency. The op-amp needs rail-to-rail inputs and outputs because it
runs as a follower from a single \qty{3.3}{\volt} rail, and a gain-bandwidth
of at least a hundred times the corner frequency keeps its own roll-off from
affecting the response~\cite{karki_sallen_key_2002}. A general-purpose CMOS
part such as the MCP6002 meets both points.
